Une tasse en céramique (émissivité \num{\v{em}}) contient \qty{250}{g} de café à \qty{\v{T}}{\degreeCelsius}. Sa surface extérieure, à la température du café, a une aire de \qty{\v{S}}{cm^2}. La pièce est à \qty{\v{Te}}{\degreeCelsius}.

\begin{enumerate}[label=\alph*)]
    \item Déterminez la puissance nette que la tasse perd par rayonnement.
    \item En supposant cette puissance constante, quelle énergie le café perd-il par rayonnement en \qty{\v{t}}{min} ?
    \item De combien de degrés le café se refroidirait-il si cette énergie était le seul mécanisme de perte ($c_{\text{eau}} = \qty{4186}{J/(kg.K)}$, on néglige la tasse) ? Quels autres mécanismes refroidissent le café ?
\end{enumerate}
